\documentclass[12pt,a4paper]{article}

% ====== ENCODING & FONT ======
\usepackage[utf8]{inputenc}    % Hỗ trợ UTF-8
\usepackage[T5]{fontenc}       % Font encoding cho tiếng Việt
\usepackage[vietnamese]{babel} % Ngôn ngữ tiếng Việt

% ====== GEOMETRY & LAYOUT ======
\usepackage{geometry}          % Tùy chỉnh lề
\geometry{
  a4paper,
  top=25mm,
  bottom=25mm,
  left=25mm,
  right=25mm
}

\usepackage{setspace}          % Chỉnh khoảng cách dòng
\onehalfspacing                % Giãn dòng 1.5

\usepackage{parskip}           % Tự động thêm khoảng cách giữa các đoạn
\setlength{\parindent}{0pt}    % Không thụt đầu dòng

% ====== MATHEMATICS ======
\usepackage{amsmath, amssymb, amsfonts} % Toán học nâng cao
\usepackage{mathtools}                  % Bổ sung cho amsmath

% ====== GRAPHICS & FIGURES ======
\usepackage{graphicx}          % Thêm hình ảnh
\usepackage{float}             % Quản lý vị trí hình/table
\usepackage{caption}           % Tùy chỉnh caption
\usepackage{subcaption}        % Sub-figures

% ====== TABLES ======
\usepackage{array}             % Cải tiến bảng
\usepackage{booktabs}          % Dòng bảng đẹp hơn
\usepackage{multirow}          % Gộp ô bảng theo chiều dọc
\usepackage{multicol}          % Gộp cột
\usepackage{longtable}         % Bảng dài nhiều trang
\usepackage{tabularx}          % Bảng có chiều rộng tự động

% ====== COLORS & HYPERLINKS ======
\usepackage[dvipsnames]{xcolor} % Màu sắc mở rộng
\usepackage{hyperref}           % Tạo liên kết
\hypersetup{
    colorlinks=true,
    linkcolor=blue,
    urlcolor=blue,
    citecolor=blue,
    pdfborder={0 0 0}
}

% ====== LISTS ======
\usepackage{enumitem}           % Tùy chỉnh danh sách
\setlist[itemize]{noitemsep}    % Giảm khoảng cách itemize
\setlist[enumerate]{noitemsep}

% ====== HEADER & FOOTER ======
\usepackage{fancyhdr}
\pagestyle{fancy}
\fancyhf{}
\fancyhead[L]{\leftmark}  % Tiêu đề bên trái
\fancyhead[R]{\thepage}   % Trang bên phải

% ====== OTHER UTILITIES ======
\usepackage{pdfpages}          % Chèn file PDF
\usepackage{lipsum}            % Sinh text giả (demo)
\usepackage{titlesec}          % Tùy chỉnh tiêu đề
\titleformat{\section}{\large\bfseries}{\thesection}{1em}{}
\titleformat{\subsection}{\normalsize\bfseries}{\thesubsection}{1em}{}


\begin{document}

\textbf{BÀI TẬP CUỐI CHƯƠNG 4}

\textbf{Phần I}. \textbf{Trắc nghiệm}

\textbf{Câu} \textbf{1.} Với mọi góc nhọn \(\alpha\), ta có

\textbf{A.} \(\cot\left( 90^{o} - \alpha \right) = cos\alpha\).
\textbf{B.} \begin{center}

\includegraphics[keepaspectratio,height=1.8ex]{./media/image1.png}
\end{center}.

\textbf{C.} \(\cot\left( 90^{o} - \alpha \right) = sin\alpha\).
\textbf{\ul{D.}} \(\tan\left( 90^{o} - \alpha \right) = cot\alpha\).

\begin{longtable}[]{@{}
  >{\raggedright\arraybackslash}p{(\columnwidth - 2\tabcolsep) * \real{0.6286}}
  >{\raggedright\arraybackslash}p{(\columnwidth - 2\tabcolsep) * \real{0.3714}}@{}}
\toprule\noalign{}
\begin{minipage}[b]{\linewidth}\raggedright
\textbf{Câu 2.} Hãy tính chiều cao của tháp Eiffel mà không cần lên tận
đỉnh tháp khi biết góc tạo bởi tia nắng mặt trời và mặt đất là
\(62^{0}\) và bóng tháp trên mặt đất khi đó là \(172\ m\) (làm tròn kết
quả tới chữ số thập phân thứ nhất)

\textbf{A.} \(323,4\ m\) .

\textbf{\ul{B.}} \(323,5\ m\).

\textbf{C.} \(323,6\ m\).

\textbf{D.} \(323,7\ m\).
\end{minipage} & \begin{minipage}[b]{\linewidth}\raggedright
\begin{center}
\includegraphics[width=1.55328in,height=1.83962in]{./media/image2.emf}
\end{center}
\end{minipage} \\ \hline\midrule\noalign{}
\endhead
\bottomrule\noalign{}
\endlastfoot
\end{longtable}

\textbf{Phần III}. \textbf{Tự luận}

\begin{longtable}[]{@{}
  >{\raggedright\arraybackslash}p{(\columnwidth - 2\tabcolsep) * \real{0.4632}}
  >{\raggedright\arraybackslash}p{(\columnwidth - 2\tabcolsep) * \real{0.5368}}@{}}
\toprule\noalign{}
\begin{minipage}[b]{\linewidth}\raggedright
\textbf{Bài 1.} Một chiếc máy bay cất cánh theo một góc \(25{^\circ}\)
so với mặt đất. Hỏi muốn đạt độ cao \(2000m\) thì máy bay phải bay một
đoạn đường là bao nhiêu mét \emph{( làm tròn kết quả đến chữ số thập
phân thứ nhất)}
\end{minipage} & \begin{minipage}[b]{\linewidth}\raggedright
\begin{center}
\includegraphics[width=3.31458in,height=1.62839in]{./media/image3.emf}
\end{center}
\end{minipage} \\ \hline\midrule\noalign{}
\endhead
\bottomrule\noalign{}
\endlastfoot
\end{longtable}

\textbf{Bài 2.} Cho \(\Delta ABC\)vuông tại \(A\) có đường cao
\(AH(H \in BC)\). Biết\(BC = 20cm\) và~\(\sin C = 0,6\). Tính độ dài
cạnh\(AB\), \(AH\), số đo góc \(B\) (\emph{số đo góc làm tròn đến độ}).


\end{document}
